\documentclass[10pt,a4paper]{amsart}
\usepackage{amsmath,amssymb,amsthm,graphicx}
% Retained preamble for the Stacks Project chapter snapshot in this directory.
% The chapter's own preamble.tex is a copy of the shared Stacks Project preamble; the block
% below is sliced MECHANICALLY, line for line, from it. Nothing was added, removed or
% redefined; the wrapper class, the xr-hyper/hyperref wiring and the externaldocument list
% are omitted because the assembler supplies the document class.
%
% lines taken from preamble.tex: 3, 5, 9, 11, 12, 156, 157, 158, ... (51 lines total)
%
% package coverage restored from the chapter's own preamble (2026-10-13): xy, xy-2cell

\usepackage[all]{xy}
\xyoption{2cell}
\UseAllTwocells

\newenvironment{reference}{\comment}{\endcomment}
\newenvironment{slogan}{\comment}{\endcomment}
\newenvironment{history}{\comment}{\endcomment}
\theoremstyle{plain}
\newtheorem{theorem}[subsection]{Theorem}
\newtheorem{proposition}[subsection]{Proposition}
\newtheorem{lemma}[subsection]{Lemma}
\theoremstyle{definition}
\newtheorem{definition}[subsection]{Definition}
\newtheorem{example}[subsection]{Example}
\newtheorem{exercise}[subsection]{Exercise}
\newtheorem{situation}[subsection]{Situation}
\theoremstyle{remark}
\newtheorem{remark}[subsection]{Remark}
\newtheorem{remarks}[subsection]{Remarks}
\numberwithin{equation}{subsection}
\def\lim{\mathop{\mathrm{lim}}\nolimits}
\def\colim{\mathop{\mathrm{colim}}\nolimits}
\def\Spec{\mathop{\mathrm{Spec}}}
\def\Hom{\mathop{\mathrm{Hom}}\nolimits}
\def\Ext{\mathop{\mathrm{Ext}}\nolimits}
\def\SheafHom{\mathop{\mathcal{H}\!\mathit{om}}\nolimits}
\def\SheafExt{\mathop{\mathcal{E}\!\mathit{xt}}\nolimits}
\def\Sch{\mathit{Sch}}
\def\Mor{\mathop{\mathrm{Mor}}\nolimits}
\def\Ob{\mathop{\mathrm{Ob}}\nolimits}
\def\Sh{\mathop{\mathit{Sh}}\nolimits}
\def\NL{\mathop{N\!L}\nolimits}
\def\CH{\mathop{\mathrm{CH}}\nolimits}
\def\proetale{{pro\text{-}\acute{e}tale}}
\def\etale{{\acute{e}tale}}
\def\QCoh{\mathit{QCoh}}
\def\Ker{\mathop{\mathrm{Ker}}}
\def\Im{\mathop{\mathrm{Im}}}
\def\Coker{\mathop{\mathrm{Coker}}}
\def\Coim{\mathop{\mathrm{Coim}}}
\DeclareMathSymbol{\boxtimes}{\mathbin}{AMSa}{"02}
\def\QCohstack{\mathcal{QC}\!\mathit{oh}}
\def\Cohstack{\mathcal{C}\!\mathit{oh}}
\def\Spacesstack{\mathcal{S}\!\mathit{paces}}
\def\Quotfunctor{\mathrm{Quot}}
\def\Hilbfunctor{\mathrm{Hilb}}
\def\Curvesstack{\mathcal{C}\!\mathit{urves}}
\def\Polarizedstack{\mathcal{P}\!\mathit{olarized}}
\def\Complexesstack{\mathcal{C}\!\mathit{omplexes}}
\def\Pic{\mathop{\mathrm{Pic}}\nolimits}
\def\Picardstack{\mathcal{P}\!\mathit{ic}}
\def\Picardfunctor{\mathrm{Pic}}
\def\Deformationcategory{\mathcal{D}\!\mathit{ef}}
\usepackage{cleveref}
\crefname{originallemma}{original lemma}{original lemmas}
\Crefname{originallemma}{Original lemma}{Original lemmas}
\makeatletter
\newlabel{lemma-Isom-in-Mor}{{0DPP}{source}{}{sourcepointer-1}{}}
\newlabel{lemma-Isom-in-Mor@cref}{{[originallemma][0][]0DPP}{[1][0][]source}{}{}{}}
\makeatother
\crefname{originallemma}{original lemma}{original lemmas}
\Crefname{originallemma}{Original lemma}{Original lemmas}
\makeatletter
\newlabel{lemma-Mor-qc-over-base}{{0DPR}{source}{}{sourcepointer-2}{}}
\newlabel{lemma-Mor-qc-over-base@cref}{{[originallemma][0][]0DPR}{[1][0][]source}{}{}{}}
\makeatother
\crefname{originallemma}{original lemma}{original lemmas}
\Crefname{originallemma}{Original lemma}{Original lemmas}
\makeatletter
\newlabel{lemma-Mor-s-lfp}{{0DPN}{source}{}{sourcepointer-3}{}}
\newlabel{lemma-Mor-s-lfp@cref}{{[originallemma][0][]0DPN}{[1][0][]source}{}{}{}}
\makeatother
\crefname{originallemma}{original lemma}{original lemmas}
\Crefname{originallemma}{Original lemma}{Original lemmas}
\makeatletter
\newlabel{perfect-lemma-chi-locally-constant-geometric}{{0B9T}{source}{}{sourcepointer-4}{}}
\newlabel{perfect-lemma-chi-locally-constant-geometric@cref}{{[originallemma][0][]0B9T}{[1][0][]source}{}{}{}}
\makeatother
\crefname{originallemma}{original lemma}{original lemmas}
\Crefname{originallemma}{Original lemma}{Original lemmas}
\makeatletter
\newlabel{quot-lemma-polarized-limits}{{0D43}{source}{}{sourcepointer-5}{}}
\newlabel{quot-lemma-polarized-limits@cref}{{[originallemma][0][]0D43}{[1][0][]source}{}{}{}}
\makeatother
\crefname{originallemma}{original lemma}{original lemmas}
\Crefname{originallemma}{Original lemma}{Original lemmas}
\makeatletter
\newlabel{remark-Mor-numerical}{{0DPQ}{source}{}{sourcepointer-6}{}}
\newlabel{remark-Mor-numerical@cref}{{[originallemma][0][]0DPQ}{[1][0][]source}{}{}{}}
\makeatother
\crefname{originallemma}{original lemma}{original lemmas}
\Crefname{originallemma}{Original lemma}{Original lemmas}
\makeatletter
\newlabel{section-hom-isom}{{0DLW}{source}{}{sourcepointer-7}{}}
\newlabel{section-hom-isom@cref}{{[originallemma][0][]0DLW}{[1][0][]source}{}{}{}}
\makeatother
\crefname{originallemma}{original lemma}{original lemmas}
\Crefname{originallemma}{Original lemma}{Original lemmas}
\makeatletter
\newlabel{stacks-limits-lemma-limit-preserving-diagonal}{{0CMW}{source}{}{sourcepointer-8}{}}
\newlabel{stacks-limits-lemma-limit-preserving-diagonal@cref}{{[originallemma][0][]0CMW}{[1][0][]source}{}{}{}}
\makeatother
\crefname{originallemma}{original lemma}{original lemmas}
\Crefname{originallemma}{Original lemma}{Original lemmas}
\makeatletter
\newlabel{stacks-limits-proposition-characterize-locally-finite-presentation}{{0CMY}{source}{}{sourcepointer-9}{}}
\newlabel{stacks-limits-proposition-characterize-locally-finite-presentation@cref}{{[originallemma][0][]0CMY}{[1][0][]source}{}{}{}}
\makeatother
\crefname{originallemma}{original lemma}{original lemmas}
\Crefname{originallemma}{Original lemma}{Original lemmas}
\makeatletter
\newlabel{stacks-morphisms-lemma-composition-finite-presentation}{{06Q3}{source}{}{sourcepointer-10}{}}
\newlabel{stacks-morphisms-lemma-composition-finite-presentation@cref}{{[originallemma][0][]06Q3}{[1][0][]source}{}{}{}}
\makeatother
\crefname{originallemma}{original lemma}{original lemmas}
\Crefname{originallemma}{Original lemma}{Original lemmas}
\makeatletter
\newlabel{stacks-morphisms-lemma-smooth-locally-finite-presentation}{{0DNP}{source}{}{sourcepointer-11}{}}
\newlabel{stacks-morphisms-lemma-smooth-locally-finite-presentation@cref}{{[originallemma][0][]0DNP}{[1][0][]source}{}{}{}}
\makeatother
\crefname{originallemma}{original lemma}{original lemmas}
\Crefname{originallemma}{Original lemma}{Original lemmas}
\makeatletter
\newlabel{stacks-properties-lemma-check-property-covering}{{04XD}{source}{}{sourcepointer-12}{}}
\newlabel{stacks-properties-lemma-check-property-covering@cref}{{[originallemma][0][]04XD}{[1][0][]source}{}{}{}}
\makeatother
\crefname{originallemma}{original lemma}{original lemmas}
\Crefname{originallemma}{Original lemma}{Original lemmas}
\makeatletter
\newlabel{varieties-lemma-numerical-polynomial-from-euler}{{0BEM}{source}{}{sourcepointer-13}{}}
\newlabel{varieties-lemma-numerical-polynomial-from-euler@cref}{{[originallemma][0][]0BEM}{[1][0][]source}{}{}{}}
\makeatother
\title{The diagonal of the polarized stack is separated and of finite presentation}
\author{The Stacks Project Authors}
\date{stacks-project:moduli.tex:local-snapshot-20261012; GNU FDL 1.2 or later}
\begin{document}
\maketitle
\noindent\textit{Editorial note.} Statement span is lines 1982-1986 of sources/stacks-moduli/pinned-original-source.tex; proof span is lines 1988-2109, the complete author proof (120 body lines).\par\medskip
\noindent\textit{Editorial note.} Navigation stubs: 13 cited-but-not-included labels are supplied as non-load-bearing pointers with their pinned tags where the snapshot tag index has them.\par\medskip
\noindent\textit{Editorial note.} extra\_packages is empty on purpose: the retained preamble already loads xy with options and a second declaration clashes in the assembler wrapper.\par\medskip
\noindent\textit{Editorial note.} Text is retained verbatim, including the author's own typos and placeholders.\par\medskip
\noindent\textit{Editorial note.} Queue index 0 of queue\_moduli.json (pb=4917, stmt 1982-1986, proof 1988-2109).\par\medskip
\noindent\textit{Editorial note.} stacks-moduli wave-1: fresh adjudication accept (2026-10-14); pool per Chief order 08:5x (dpa superseded to lane-s9); moduli.tex 53 labels/48 with proof judged 25 accept / 23 reject / 5 pb=0 env-dup rows excluded; spans cite-free and reference/slogan/history-free so no patchenv needed.\par\medskip
\noindent\textit{Editorial note.} Cross-chapter same-label warning (NOT an occupancy conflict, Chief ruling 2026-10-13 evening): label 'lemma-polarized-diagonal-separated-fp' also exists in stacks-snapshot-20261011 — a different proposition in a different chapter file.\par\medskip
\noindent\textit{Allowed standard prerequisites (editorial).}\begin{itemize}
\item Explicit prerequisite boundary. The cited inputs are prerequisites; the argument in the retained span is not replaced by any cited result.
\end{itemize}\medskip
\ifdefined\hypertarget\hypertarget{sourcepointer-1}{}\fi
\noindent\textit{Source reference.} Original lemma 0DPP: sources/stacks-moduli/pinned-original-source.tex, source lines 1857--1862. Cited by the retained statement or proof but not part of the collected spans; supplied only as a navigation pointer. The argument does not depend on its content.\par\medskip
\ifdefined\hypertarget\hypertarget{sourcepointer-2}{}\fi
\noindent\textit{Source reference.} Original lemma 0DPR: sources/stacks-moduli/pinned-original-source.tex, source lines 1899--1911. Cited by the retained statement or proof but not part of the collected spans; supplied only as a navigation pointer. The argument does not depend on its content.\par\medskip
\ifdefined\hypertarget\hypertarget{sourcepointer-3}{}\fi
\noindent\textit{Source reference.} Original lemma 0DPN: sources/stacks-moduli/pinned-original-source.tex, source lines 1843--1846. Cited by the retained statement or proof but not part of the collected spans; supplied only as a navigation pointer. The argument does not depend on its content.\par\medskip
\ifdefined\hypertarget\hypertarget{sourcepointer-4}{}\fi
\noindent\textit{Source reference.} Original lemma 0B9T: sources/stacks-moduli/pinned-original-source.tex, source lines 2109--2109. Cited by the retained statement or proof but not part of the collected spans; supplied only as a navigation pointer. The argument does not depend on its content.\par\medskip
\ifdefined\hypertarget\hypertarget{sourcepointer-5}{}\fi
\noindent\textit{Source reference.} Original lemma 0D43: sources/stacks-moduli/pinned-original-source.tex, source lines 2109--2109. Cited by the retained statement or proof but not part of the collected spans; supplied only as a navigation pointer. The argument does not depend on its content.\par\medskip
\ifdefined\hypertarget\hypertarget{sourcepointer-6}{}\fi
\noindent\textit{Source reference.} Original lemma 0DPQ: sources/stacks-moduli/pinned-original-source.tex, source lines 1870--1896. Cited by the retained statement or proof but not part of the collected spans; supplied only as a navigation pointer. The argument does not depend on its content.\par\medskip
\ifdefined\hypertarget\hypertarget{sourcepointer-7}{}\fi
\noindent\textit{Source reference.} Original lemma 0DLW: sources/stacks-moduli/pinned-original-source.tex, source lines 296--341. Cited by the retained statement or proof but not part of the collected spans; supplied only as a navigation pointer. The argument does not depend on its content.\par\medskip
\ifdefined\hypertarget\hypertarget{sourcepointer-8}{}\fi
\noindent\textit{Source reference.} Original lemma 0CMW: sources/stacks-moduli/pinned-original-source.tex, source lines 2109--2109. Cited by the retained statement or proof but not part of the collected spans; supplied only as a navigation pointer. The argument does not depend on its content.\par\medskip
\ifdefined\hypertarget\hypertarget{sourcepointer-9}{}\fi
\noindent\textit{Source reference.} Original lemma 0CMY: sources/stacks-moduli/pinned-original-source.tex, source lines 2109--2109. Cited by the retained statement or proof but not part of the collected spans; supplied only as a navigation pointer. The argument does not depend on its content.\par\medskip
\ifdefined\hypertarget\hypertarget{sourcepointer-10}{}\fi
\noindent\textit{Source reference.} Original lemma 06Q3: sources/stacks-moduli/pinned-original-source.tex, source lines 2109--2109. Cited by the retained statement or proof but not part of the collected spans; supplied only as a navigation pointer. The argument does not depend on its content.\par\medskip
\ifdefined\hypertarget\hypertarget{sourcepointer-11}{}\fi
\noindent\textit{Source reference.} Original lemma 0DNP: sources/stacks-moduli/pinned-original-source.tex, source lines 2109--2109. Cited by the retained statement or proof but not part of the collected spans; supplied only as a navigation pointer. The argument does not depend on its content.\par\medskip
\ifdefined\hypertarget\hypertarget{sourcepointer-12}{}\fi
\noindent\textit{Source reference.} Original lemma 04XD: sources/stacks-moduli/pinned-original-source.tex, source lines 2109--2109. Cited by the retained statement or proof but not part of the collected spans; supplied only as a navigation pointer. The argument does not depend on its content.\par\medskip
\ifdefined\hypertarget\hypertarget{sourcepointer-13}{}\fi
\noindent\textit{Source reference.} Original lemma 0BEM: sources/stacks-moduli/pinned-original-source.tex, source lines 2109--2109. Cited by the retained statement or proof but not part of the collected spans; supplied only as a navigation pointer. The argument does not depend on its content.\par\medskip
\begin{lemma}
\label{lemma-polarized-diagonal-separated-fp}
The diagonal of $\Polarizedstack$ is separated
and of finite presentation.
\end{lemma}

\begin{proof}
Recall that $\Polarizedstack$ is a limit preserving algebraic stack, see
Quot, Lemma \ref{quot-lemma-polarized-limits}.
By Limits of Stacks, Lemma \ref{stacks-limits-lemma-limit-preserving-diagonal}
this implies that
$\Delta : \Polarizedstack \to \Polarizedstack \times \Polarizedstack$
is limit preserving. Hence $\Delta$ is locally of finite presentation
by Limits of Stacks, Proposition
\ref{stacks-limits-proposition-characterize-locally-finite-presentation}.

\medskip\noindent
Let us prove that $\Delta$ is separated. To see this, it suffices to show
that given an affine scheme $U$ and two objects
$\upsilon = (Y, \mathcal{N})$ and $\chi = (X, \mathcal{L})$
of $\Polarizedstack$ over $U$, the algebraic
space
$$
\mathit{Isom}_{\Polarizedstack}(\upsilon, \chi)
$$
is separated. The rule which to an isomorphism $\upsilon_T \to \chi_T$
assigns the underlying isomorphism $Y_T \to X_T$ defines a morphism
$$
\mathit{Isom}_{\Polarizedstack}(\upsilon, \chi)
\longrightarrow
\mathit{Isom}_U(Y, X)
$$
Since we have seen in Lemmas \ref{lemma-Mor-s-lfp} and
\ref{lemma-Isom-in-Mor} that the target is
a separated algebraic space, it suffices to prove that this morphism
is separated. Given an isomorphism $f : Y_T \to X_T$
over some scheme $T/U$, then clearly
$$
\mathit{Isom}_{\Polarizedstack}(\upsilon, \chi)
\times_{\mathit{Isom}_U(Y, X), [f]} T
=
\mathit{Isom}(\mathcal{N}_T, f^*\mathcal{L}_T)
$$
Here $[f] : T \to \mathit{Isom}_U(Y, X)$ indicates the $T$-valued
point corresponding to $f$ and
$\mathit{Isom}(\mathcal{N}_T, f^*\mathcal{L}_T)$
is the algebraic space discussed in Section \ref{section-hom-isom}.
Since this algebraic space is affine over $U$, the claim implies
$\Delta$ is separated.

\medskip\noindent
To finish the proof we show that $\Delta$ is quasi-compact. Since
$\Delta$ is representable by algebraic spaces, it suffice to check
the base change of $\Delta$ by a surjective smooth morphism
$U \to \Polarizedstack \times \Polarizedstack$ is quasi-compact
(see for example Properties of Stacks, Lemma
\ref{stacks-properties-lemma-check-property-covering}).
We can assume $U = \coprod U_i$ is a disjoint union of affine opens.
Since $\Polarizedstack$ is limit preserving (see above), we
see that $\Polarizedstack \to \Spec(\mathbf{Z})$ is locally of
finite presentation, hence $U_i \to \Spec(\mathbf{Z})$ is
locally of finite presentation
(Limits of Stacks, Proposition
\ref{stacks-limits-proposition-characterize-locally-finite-presentation}
and Morphisms of Stacks, Lemmas
\ref{stacks-morphisms-lemma-composition-finite-presentation} and
\ref{stacks-morphisms-lemma-smooth-locally-finite-presentation}).
In particular, $U_i$ is Noetherian affine. This reduces us to the
case discussed in the next paragraph.

\medskip\noindent
In this paragraph, given a Noetherian affine scheme $U$ and two objects
$\upsilon = (Y, \mathcal{N})$ and $\chi = (X, \mathcal{L})$
of $\Polarizedstack$ over $U$, we show the algebraic space
$$
\mathit{Isom}_{\Polarizedstack}(\upsilon, \chi)
$$
is quasi-compact. Since the connected components of $U$ are open and closed
we may replace $U$ by these. Thus we may and do assume $U$ is connected.
Let $u \in U$ be a point. Let $P$ be the Hilbert polynomial
$n \mapsto \chi(Y_u, \mathcal{N}_u^{\otimes n})$, see
Varieties, Lemma \ref{varieties-lemma-numerical-polynomial-from-euler}.
Since $U$ is connected and since
the functions $u \mapsto \chi(Y_u, \mathcal{N}_u^{\otimes n})$
are locally constant (see 
Derived Categories of Schemes, Lemma
\ref{perfect-lemma-chi-locally-constant-geometric})
we see that we get the same Hilbert polynomial in every point of $U$.
Set
$\mathcal{M} = \text{pr}_1^*\mathcal{N}
\otimes_{\mathcal{O}_{Y \times_U X}} \text{pr}_2^*\mathcal{L}$
on $Y \times_U X$. Given
$(f, \varphi) \in \mathit{Isom}_{\Polarizedstack}(\upsilon, \chi)(T)$
for some scheme $T$ over $U$ then for every $t \in T$ we have
$$
\chi(Y_t, (\text{id} \times f)^*\mathcal{M}^{\otimes n}) =
\chi(Y_t,
\mathcal{N}_t^{\otimes n} \otimes_{\mathcal{O}_{Y_t}}
f_t^*\mathcal{L}_t^{\otimes n}) =
\chi(Y_t, \mathcal{N}_t^{\otimes 2n}) = P(2n)
$$
where in the middle equality we use the isomorphism
$\varphi : f^*\mathcal{L}_T \to \mathcal{N}_T$.
Setting $P'(t) = P(2t)$ we find that the morphism
$$
\mathit{Isom}_{\Polarizedstack}(\upsilon, \chi)
\longrightarrow
\mathit{Isom}_U(Y, X)
$$
(see earlier) has image contained in the intersection
$$
\mathit{Isom}_U(Y, X) \cap \mathit{Mor}^{P', \mathcal{M}}_U(Y, X)
$$
The intersection is an intersection of open subspaces of
$\mathit{Mor}_U(Y, X)$ (see Lemma \ref{lemma-Isom-in-Mor} and
Remark \ref{remark-Mor-numerical}).
Now $\mathit{Mor}^{P', \mathcal{M}}_U(Y, X)$
is a Noetherian algebraic space as it is of finite
presentation over $U$ by Lemma \ref{lemma-Mor-qc-over-base}.
Thus the intersection
is a Noetherian algebraic space too. Since the morphism
$$
\mathit{Isom}_{\Polarizedstack}(\upsilon, \chi)
\longrightarrow
\mathit{Isom}_U(Y, X) \cap \mathit{Mor}^{P', \mathcal{M}}_U(Y, X)
$$
is affine (see above) we conclude.
\end{proof}

\end{document}
